\documentclass[12pt,a4paper,oneside]{report}

% ==============================================================================
% 1. CẤU HÌNH NGÔN NGỮ TIẾNG VIỆT & PHÔNG CHỮ (CHUẨN OVERLEAF / PDFLATEX)
% ==============================================================================
\usepackage[utf8]{inputenc}
\usepackage[vietnamese]{babel}
\usepackage[T5]{fontenc}
\usepackage{lmodern}

% ==============================================================================
% 2. CĂN CHỈNH TRANG IN (CANH LỀ CHUẨN ĐẠI HỌC FPT)
% Lề trên: 2.5cm, Lề dưới: 2.5cm, Lề trái (đóng gáy): 3.0cm, Lề phải: 2.0cm
% ==============================================================================
\usepackage[a4paper,top=2.5cm,bottom=2.5cm,left=3.0cm,right=2.0cm]{geometry}
\usepackage{setspace}
\onehalfspacing % Giãn dòng 1.5 lines tiêu chuẩn học thuật

% ==============================================================================
% 3. CÁC GÓI BỔ TRỢ ĐỒ HỌA, MÀU SẮC & BẢNG BIỂU
% ==============================================================================
\usepackage{graphicx}
\usepackage[table]{xcolor}
\usepackage{booktabs}
\usepackage{tabularx}
\usepackage{multirow}
\usepackage{amsmath,amssymb}
\usepackage{tikz}
\usetikzlibrary{shapes.geometric, arrows, positioning, fit, calc}
\usepackage[many]{tcolorbox}
\usepackage{caption}
\usepackage{subcaption}

% ==============================================================================
% 4. ĐỊNH DẠNG MÃ NGUỒN (CODE LISTINGS)
% ==============================================================================
\usepackage{listings}
\definecolor{codebg}{rgb}{0.97,0.98,0.99}
\definecolor{codeblue}{rgb}{0.12,0.47,0.71}
\definecolor{codegreen}{rgb}{0.18,0.55,0.34}
\definecolor{codegray}{rgb}{0.5,0.5,0.5}
\definecolor{codered}{rgb}{0.75,0.22,0.17}

\lstdefinestyle{dmsStyle}{
    backgroundcolor=\color{codebg},   
    commentstyle=\color{codegreen}\itshape,
    keywordstyle=\color{codeblue}\bfseries,
    numberstyle=\tiny\color{codegray},
    stringstyle=\color{codered},
    basicstyle=\ttfamily\footnotesize,
    breakatwhitespace=false,         
    breaklines=true,                 
    captionpos=b,                    
    keepspaces=true,                 
    numbers=left,                    
    numbersep=7pt,                  
    showspaces=false,                
    showstringspaces=false,
    showtabs=false,                  
    tabsize=2,
    frame=single,
    rulecolor=\color{gray!30}
}
\lstset{style=dmsStyle}

% ==============================================================================
% 5. CẤU HÌNH LIÊN KẾT ĐIỆN TỬ (HYPERREF & BOOKMARKS)
% ==============================================================================
\usepackage{hyperref}
\hypersetup{
    colorlinks=true,
    linkcolor=blue!75!black,
    citecolor=blue!75!black,
    urlcolor=teal!85!black,
    pdfauthor={FER202_Group4},
    pdftitle={Báo cáo Dự án Domanage - Hệ thống Quản lý KTX Trực tuyến},
    pdfsubject={Front-End Web Development with React},
    pdfkeywords={React 19, TypeScript, Tailwind CSS, Dormitory Management, Overleaf}
}

% ==============================================================================
% 6. HEADER & FOOTER TIÊU CHUẨN (FANCYHDR)
% ==============================================================================
\usepackage{fancyhdr}
\pagestyle{fancy}
\fancyhf{}
\fancyhead[L]{\slshape\nouppercase{\leftmark}}
\fancyhead[R]{\slshape FER202 -- Nhóm 4}
\fancyfoot[C]{\thepage}
\renewcommand{\headrulewidth}{0.4pt}
\renewcommand{\footrulewidth}{0.4pt}

% ==============================================================================
% BẮT ĐẦU TÀI LIỆU
% ==============================================================================
\begin{document}

% Trang bìa
\begin{titlepage}
\begin{center}

% Header Đại học FPT
\vspace*{-1cm}
{\fontsize{14pt}{16pt}\selectfont \textbf{BỘ GIÁO DỤC VÀ ĐÀO TẠO}}\\[0.2cm]
{\fontsize{16pt}{18pt}\selectfont \textbf{TRƯỜNG ĐẠI HỌC FPT}}\\[0.2cm]
{\fontsize{13pt}{15pt}\selectfont \textbf{KHOA CÔNG NGHỆ THÔNG TIN}}\\[0.3cm]
\rule{0.6\textwidth}{1pt}\\[1.5cm]

% Logo placeholder (dùng TikZ nếu không có file ảnh logo)
\begin{tikzpicture}
    \draw[fill=orange!10, draw=orange!80, line width=1.5pt, rounded corners=10pt] (-3.5,-1) rectangle (3.5,1);
    \node at (0,0.3) {\fontsize{16pt}{18pt}\selectfont \textbf{\color{orange!90!black}FPT UNIVERSITY}};
    \node at (0,-0.3) {\fontsize{10pt}{12pt}\selectfont \textit{Khát Vọng Đổi Thay - Leading Technology}};
\end{tikzpicture}\\[1.5cm]

% Tên môn học
{\fontsize{14pt}{16pt}\selectfont \textbf{BÁO CÁO DỰ ÁN MÔN HỌC}}\\[0.2cm]
{\fontsize{15pt}{18pt}\selectfont \textbf{MÔN HỌC: FER202 -- FRONT-END WEB DEVELOPMENT WITH REACT}}\\[1.0cm]

% Tên Đề tài
\begin{tcolorbox}[colback=blue!5!white, colframe=blue!75!black, arc=6mm, width=0.95\textwidth, halign=center, boxsep=3mm]
    {\fontsize{17pt}{22pt}\selectfont \textbf{\color{blue!80!black}ĐỀ TÀI:}}\\[0.2cm]
    {\fontsize{20pt}{26pt}\selectfont \textbf{\color{blue!90!black}DOMANAGE -- HỆ THỐNG QUẢN LÝ KÝ TÚC XÁ TRỰC TUYẾN}}\\[0.2cm]
    {\fontsize{13pt}{16pt}\selectfont \textit{(Online Dormitory Management System -- DMS)}}
\end{tcolorbox}\\[1.2cm]

% Bảng thông tin sinh viên & Giảng viên
\vfill
\begin{minipage}{0.9\textwidth}
\begin{flushleft}
\fontsize{12pt}{15pt}\selectfont
\textbf{Giảng viên hướng dẫn:} ThS. Giảng viên bộ môn FER202\\[0.3cm]
\textbf{Nhóm thực hiện:} FER202\_Group4\\[0.3cm]
\textbf{Danh sách thành viên nhóm:}\\[0.15cm]
\begin{tabular}{|c|l|c|l|}
\hline
\textbf{STT} & \textbf{Họ và Tên} & \textbf{MSSV} & \textbf{Vai trò chính trong dự án} \\ \hline
1 & Nguyễn Hoàng Vinh & SE170001 & Project Manager \& Full-Stack Lead \\ \hline
2 & Trần Công Tiến & SE170002 & Frontend Lead \& UI/UX Designer \\ \hline
3 & Trần Thanh Nguyên & SE170003 & Backend Engineer \& Database Architect \\ \hline
4 & Văn Thái Trung & SE170004 & Frontend Developer \& Integration Specialist \\ \hline
5 & Đinh Thanh Trung & SE170005 & Backend Developer \& Security Analyst \\ \hline
6 & Huỳnh Gia Kim & SE170006 & Quality Assurance (QA) \& Test Engineer \\ \hline
\end{tabular}
\end{flushleft}
\end{minipage}

\vfill
{\fontsize{12pt}{14pt}\selectfont \textbf{TP. HỒ CHÍ MINH, THÁNG 10/2026}}

\end{center}
\end{titlepage}



% Lời cam đoan
\newpage
\pagenumbering{roman}
\chapter*{LỜI CAM ĐOAN}
Chúng tôi xin cam đoan đề tài nghiên cứu và xây dựng hệ thống \textbf{``Domanage -- Hệ thống Quản lý Ký túc xá Trực tuyến (Online Dormitory Management System)''} là sản phẩm nghiên cứu và phát triển độc lập của tập thể nhóm 4 môn học FER202 (Front-End Web Development with React). 

Tất cả các số liệu, sơ đồ, kiến trúc và mã nguồn trình bày trong báo cáo này hoàn toàn trung thực, phản ánh chính xác kết quả làm việc thực tế của nhóm; không sao chép bất hợp pháp từ bất kỳ công trình hay đồ án nào khác. Mọi tài liệu tham khảo và thư viện nguồn mở sử dụng trong dự án đều được trích dẫn nguồn đầy đủ theo tiêu chuẩn học thuật.

\vspace{1.5cm}
\begin{flushright}
\textit{TP. Hồ Chí Minh, ngày 29 tháng 10 năm 2026}\\[0.3cm]
\textbf{Đại diện Nhóm thực hiện -- Nhóm trưởng}\\[1.5cm]
\textbf{Nguyễn Hoàng Vinh}
\end{flushright}

% Lời cảm ơn
\newpage
\chapter*{LỜI CẢM ƠN}
Để hoàn thành xuất sắc dự án và bản báo cáo này, tập thể Nhóm 4 xin gửi lời tri ân sâu sắc nhất tới:
\begin{itemize}
    \item \textbf{Ban Giám hiệu \& Khoa Công nghệ Thông tin -- Trường Đại học FPT:} Đã tạo dựng môi trường học thuật chuyên nghiệp, giàu tính thực tiễn cùng cơ sở vật chất hiện đại hỗ trợ sinh viên trong suốt kỳ học.
    \item \textbf{Giảng viên bộ môn FER202:} Thầy/Cô đã tận tình truyền đạt những kiến thức quý báu về tư duy lập trình giao diện người dùng hiện đại với React, các best practices trong kiến trúc Front-end, và luôn theo sát, góp ý định hướng quý báu cho đề tài của nhóm từ những ngày đầu.
    \item \textbf{Toàn thể các thành viên Nhóm 4:} Đã luôn giữ vững tinh thần trách nhiệm, chủ động học hỏi, đoàn kết và nỗ lực hết mình để đưa dự án cán đích thành công vượt tiến độ.
\end{itemize}

Nhóm rất mong nhận được những ý kiến đóng góp quý báu từ Hội đồng Giảng viên để hệ thống ngày càng hoàn thiện hơn nữa!

\vspace{1.5cm}
\begin{flushright}
\textbf{Tập thể sinh viên Nhóm 4 -- FER202}
\end{flushright}

% Tóm tắt đề tài (Abstract)
\newpage
\chapter*{TÓM TẮT ĐỀ TÀI (ABSTRACT)}
\addcontentsline{toc}{chapter}{Tóm tắt đề tài (Abstract)}

\textbf{Tóm tắt:} Dự án \textbf{Domanage (DMS)} tập trung giải quyết các bài toán bất cập trong công tác quản trị ký túc xá đại học quy mô lớn. Ứng dụng được xây dựng trên nền tảng Single Page Application với React 19, TypeScript và Tailwind CSS v4. Đề tài đem lại 3 điểm đột phá sáng tạo: (1) \textbf{Mô hình không gian 1:4:6} trực quan hóa việc chọn giường ngủ, xóa bỏ hoàn toàn cảm giác mơ hồ khi đăng ký phòng; (2) \textbf{Động cơ Tự động xét duyệt (Auto-Approval Engine)} 3 điều kiện và tự động kết xuất Hợp đồng lưu trú PDF trong dưới 3 giây; (3) \textbf{Thuật toán phân bổ chi phí điện nước ($Total / N$)} gắn liền với hình ảnh chụp công tơ thực tế làm bằng chứng đối soát minh bạch. Kết quả kiểm thử thực nghiệm trên 28 Use Cases và 90 ca kiểm thử đạt tỷ lệ thành công 100\%, chứng minh tính khả thi và tiềm năng ứng dụng thực tế cao của đề tài.

\vspace{0.8cm}
\noindent \textbf{Từ khóa:} \textit{Ký túc xá trực tuyến, React 19, TypeScript, Tailwind CSS, Auto-Approval Engine, Split-Billing, Mô hình 1:4:6, Thẻ QR Giới nghiêm.}

\vspace{1.5cm}
\noindent \rule{\textwidth}{0.5pt}\\[0.5cm]

\textbf{Abstract:} The \textbf{Domanage (DMS)} project addresses key management and operational challenges in large-scale university dormitory accommodation. Built as a high-performance Single Page Application using React 19, TypeScript, and Tailwind CSS v4, the project delivers three primary innovations: (1) An \textbf{interactive 1:4:6 spatial topology} for bed allocation that eliminates uncertainty during room selection; (2) A 3-step \textbf{Auto-Approval Engine} that verifies student eligibility and digitally renders verifiable PDF contracts in under 3 seconds; and (3) A transparent \textbf{split-billing algorithm ($Total / N$)} supported by physical meter photo evidence. Empirical testing across 28 functional Use Cases and 90 test cases achieved a 100\% pass rate, demonstrating strong architectural soundness and high readiness for real-world deployment.

\vspace{0.8cm}
\noindent \textbf{Keywords:} \textit{Dormitory Management System, React 19, TypeScript, Tailwind CSS, Auto-Approval Engine, Split-Billing, 1:4:6 Topology, Dynamic QR Curfew.}

% Mục lục & Danh mục bảng biểu
\newpage
\tableofcontents

\newpage
\listoffigures
\addcontentsline{toc}{chapter}{\listfigurename}

\newpage
\listoftables
\addcontentsline{toc}{chapter}{\listtablename}

% ==============================================================================
% NỘI DUNG CHÍNH (ĐÁNH SỐ TRANG 1, 2, 3...)
% ==============================================================================
\newpage
\pagenumbering{arabic}

\chapter{GIỚI THIỆU VÀ ĐẶT VẤN ĐỀ}
\label{ch:introduction}

\section{Bối cảnh và Lý do chọn đề tài}
Trong môi trường giáo dục đại học hiện đại, ký túc xá (KTX) không chỉ đơn thuần là nơi lưu trú mà còn là một hệ sinh thái sinh hoạt, học tập và rèn luyện quan trọng của sinh viên. Tại các cơ sở đào tạo quy mô lớn như Đại học FPT, số lượng sinh viên nội trú lên tới hàng ngàn người mỗi học kỳ, đặt ra thách thức quản trị vô cùng lớn cho Ban Quản lý KTX.

Trước đây, phần lớn các nghiệp vụ tại ký túc xá vẫn được thực hiện theo phương thức bán thủ công hoặc dựa trên các tệp bảng tính rời rạc (Google Sheets, Excel):
\begin{itemize}
    \item \textbf{Đăng ký và xếp phòng thiếu trực quan:} Sinh viên không thể hình dung vị trí cụ thể của giường ngủ trong phòng (hướng cửa sổ, tầng trên hay tầng dưới), dẫn tới tình trạng bất mãn khi nhận phòng thực tế.
    \item \textbf{Quy trình phê duyệt hợp đồng chậm trễ:} Việc đối chiếu hồ sơ, kiểm tra điều kiện sinh viên (giới tính, công nợ kỳ trước, kỷ luật) và xác nhận tiền cọc tốn nhiều ngày làm việc của cán bộ.
    \item \textbf{Bất cập và thiếu minh bạch trong tính tiền điện nước:} Tiền điện nước phòng chung thường gây tranh cãi gay gắt giữa các bạn cùng phòng do thiếu bằng chứng chỉ số công tơ, không rõ công thức chia sẻ chi phí khi có người chuyển vào hoặc chuyển đi giữa kỳ.
    \item \textbf{Báo hỏng và bảo trì cơ sở vật chất chậm trễ:} Phiếu báo hỏng bằng giấy hoặc gửi qua mạng xã hội dễ bị thất lạc, kỹ thuật viên không nắm rõ mức độ khẩn cấp, sinh viên không theo dõi được tiến độ sửa chữa.
    \item \textbf{Kiểm soát an ninh và giờ giới nghiêm 23:00 thiếu chính xác:} Việc ghi sổ giấy tại cổng bảo vệ dễ dẫn tới sai sót, gian lận quẹt thẻ hộ, và khó đối chiếu với các đơn xin về muộn đã được duyệt trước đó.
\end{itemize}

Xuất phát từ những bất cập thực tế nêu trên, dự án \textbf{``Domanage -- Online Dormitory Management System (DMS)''} được nghiên cứu và phát triển trong khuôn khổ môn học \textit{FER202 (Front-End Web Development with React)}. Dự án hướng tới mục tiêu hiện đại hóa, tự động hóa và minh bạch hóa toàn bộ vòng đời lưu trú của sinh viên tại ký túc xá trường đại học.

\section{Mục tiêu của Dự án}
Hệ thống \textbf{Domanage} được thiết kế với các mục tiêu trọng tâm sau:
\begin{enumerate}
    \item \textbf{Trải nghiệm người dùng vượt trội (User-Centric UI/UX):} Cung cấp bản đồ chọn giường trực quan theo phân cấp không gian \textbf{1:4:6} (1 Tòa nhà $\rightarrow$ Nhiều Tầng $\rightarrow$ Các Phòng $\rightarrow$ 4 Buồng cá nhân $\rightarrow$ 6 Giường/buồng), mô phỏng chân thực layout thực tế của phòng ở KTX.
    \item \textbf{Cơ chế giữ chỗ và thanh toán tức thì:} Áp dụng cơ chế khóa giường tạm thời 15 phút kèm đồng hồ đếm ngược, tích hợp mã chuyển khoản VietQR động theo chuẩn Napas 247, giúp sinh viên hoàn tất đặt chỗ nhanh chóng.
    \item \textbf{Động cơ tự động xét duyệt (Auto-Approval Engine):} Xóa bỏ hoàn toàn khâu duyệt thủ công bằng thuật toán tự động kiểm tra 3 điều kiện (khớp giới tính tòa nhà, không nợ kỳ trước/không kỷ luật, xác nhận cọc 100\%) và tự động xuất Hợp đồng điện tử (e-Contract PDF) có mã định danh xác thực.
    \item \textbf{Minh bạch hóa thu phí điện nước phòng:} Xây dựng thuật toán phân bổ chi phí tự động ($Total / N$) kèm tính năng Pro-rata theo số ngày lưu trú, lưu trữ hình ảnh chụp thực tế công tơ điện nước làm bằng chứng đối soát minh bạch.
    \item \textbf{Vận hành khép kín quy trình bảo trì & an ninh:} Số hóa toàn diện quy trình tiếp nhận phiếu báo hỏng thiết bị kèm ảnh chụp, bảng điều khiển kỹ thuật viên (Technician Bounty Board), và hệ thống quản lý giờ giới nghiêm 23:00 với mã QR động chống gian lận và trừ điểm kỷ luật KTX.
\end{enumerate}

\section{Phạm vi của Đề tài}
Dự án tập trung vào việc nghiên cứu, thiết kế và hiện thực hóa ứng dụng Web hiện đại:
\begin{itemize}
    \item \textbf{Công nghệ cốt lõi:} Single Page Application (SPA) xây dựng trên nền tảng React 19, TypeScript, Vite, Tailwind CSS v4 và Radix UI / ShadCN UI.
    \item \textbf{Mô phỏng dữ liệu & Backend:} Sử dụng JSON-Server với Engine tự động xét duyệt tùy biến (custom endpoints), kết hợp kiến trúc RESTful API và lưu trữ trạng thái đồng bộ LocalStorage/SessionStorage.
    \item \textbf{Đa nền tảng và Responsive:} Hệ thống hỗ trợ tối ưu hiển thị trên màn hình Desktop (cho Cán bộ quản lý KTX) và giao diện Mobile Web/PWA (cho Sinh viên tra cứu hóa đơn, thẻ QR ra vào cổng và Bảo vệ trực chốt).
\end{itemize}

\section{Đối tượng sử dụng hệ thống (Stakeholders \& Actors)}
Hệ thống phân quyền truy cập nghiêm ngặt dựa trên vai trò (Role-Based Access Control -- RBAC) với 5 nhóm người dùng chính:
\begin{enumerate}
    \item \textbf{Sinh viên (Student):} Đăng nhập bằng Gmail cá nhân/trường, thực hiện đăng ký và chọn giường trực quan, thanh toán cọc, nhận hợp đồng điện tử, theo dõi hóa đơn điện nước phòng, nộp phiếu báo hỏng thiết bị, và xuất trình thẻ QR di động khi qua cổng KTX sau 23:00.
    \item \textbf{Cán bộ Quản lý KTX (Dorm Manager):} Giám sát công suất phòng và tỷ lệ lấp đầy qua biểu đồ trực quan, phát hành hóa đơn điện nước hàng tháng, phê duyệt đơn xin về muộn ngoài giờ, xử lý biên bản vi phạm kỷ luật, và xuất báo cáo lưu trú/công nợ sang định dạng Excel.
    \item \textbf{Kỹ thuật viên (Technician):} Tiếp nhận các phiếu báo hỏng thiết bị theo bảng việc (Bounty Board), ghi nhận vật tư thay thế, chụp ảnh nghiệm thu hoàn thành sửa chữa, và thực hiện ghi chỉ số công tơ điện nước định kỳ kèm ảnh chụp đồng hồ.
    \item \textbf{Nhân viên Bảo vệ (Security Guard):} Sử dụng máy quét hoặc giao diện tra cứu nhanh MSSV tại cổng KTX, đối chiếu đơn xin về muộn đã duyệt, tự động ghi nhận biên bản vi phạm giờ giới nghiêm sau 23:00.
    \item \textbf{Quản trị viên Hệ thống (System Administrator):} Quản lý danh mục cơ sở vật chất (Tòa nhà, Tầng, Phòng, Giường), cấu hình đơn giá điện nước và thuế VAT, phân quyền tài khoản người dùng và bảo trì dữ liệu hệ thống.
\end{enumerate}

\section{Bố cục của Báo cáo}
Nội dung báo cáo dự án được kết cấu thành 7 chương chính như sau:
\begin{itemize}
    \item \textbf{Chương 1: Giới thiệu và Đặt vấn đề:} Trình bày bối cảnh, tính cấp thiết, mục tiêu và phạm vi đề tài.
    \item \textbf{Chương 2: Phân tích Yêu cầu Hệ thống:} Đặc tả 28 Use Cases, yêu cầu phi chức năng (NFR) và các quy tắc nghiệp vụ (Business Rules).
    \item \textbf{Chương 3: Kiến trúc Hệ thống và Thiết kế Dữ liệu:} Kiến trúc kỹ thuật, mô hình dữ liệu 12 thực thể chuẩn 3NF, mô hình không gian 1:4:6 và luồng xử lý nghiệp vụ trọng tâm.
    \item \textbf{Chương 4: Hiện thực các Phân hệ Cốt lõi:} Chi tiết giải pháp kỹ thuật và mã nguồn tiêu biểu cho 7 phân hệ chức năng.
    \item \textbf{Chương 5: Đảm bảo Chất lượng và Kiểm thử:} Kế hoạch kiểm thử, 4 ca kiểm thử biên (Edge-cases) then chốt, ma trận kết quả kiểm thử và đánh giá độ phủ.
    \item \textbf{Chương 6: Quản trị Dự án và Phân công Công việc:} Ma trận trách nhiệm RACI, lộ trình Sprint 5 tuần và đánh giá mức độ hoàn thành của từng thành viên.
    \item \textbf{Chương 7: Kết luận và Hướng phát triển Tương lai:} Tổng kết kết quả đạt được, hạn chế còn tồn đọng và định hướng mở rộng hệ thống.
\end{itemize}


\chapter{PHÂN TÍCH YÊU CẦU HỆ THỐNG}
\label{ch:requirements}

\section{Tổng quan Yêu cầu và Danh mục Actor}
Hệ thống Ký túc xá Trực tuyến (DMS) được phân tích dựa trên tiêu chuẩn tài liệu đặc tả yêu cầu phần mềm IEEE Std 830-1998. Hệ thống phục vụ 5 tác nhân (Actors) với các quyền hạn được phân định rõ ràng như được tóm tắt trong Bảng~\ref{tab:actors}.

\begin{table}[htbp]
\centering
\caption{Danh mục các Tác nhân (Actors) trong Hệ thống Domanage}
\label{tab:actors}
\small
\begin{tabularx}{\textwidth}{|l|l|X|}
\hline
\rowcolor{blue!10}
\textbf{Tác nhân (Actor)} & \textbf{Vai trò hệ thống} & \textbf{Mô tả quyền hạn và phạm vi tác vụ} \\ \hline
\textbf{Sinh viên (Student)} & \texttt{role: student} & Đăng nhập tài khoản cá nhân, xem bản đồ giường, giữ chỗ 15 phút, thanh toán VietQR, nhận hợp đồng PDF, tra cứu hóa đơn chia đều, báo hỏng thiết bị, nộp đơn về muộn và xuất trình thẻ QR ra vào. \\ \hline
\textbf{Quản lý KTX (Manager)} & \texttt{role: manager} & Quản lý danh mục phòng, duyệt đơn về muộn, phát hành hóa đơn điện nước định kỳ, giải quyết khiếu nại, giám sát tỷ lệ lấp đầy phòng, xuất báo cáo công nợ/tạm trú Excel. \\ \hline
\textbf{Kỹ thuật viên (Technician)} & \texttt{role: technician} & Nhận phiếu sửa chữa (Bounty Board), cập nhật tiến độ, ghi nhận linh kiện thay thế, chụp ảnh hoàn thành sửa chữa, ghi chỉ số công tơ điện nước phòng ngày 25--28. \\ \hline
\textbf{Bảo vệ (Guard)} & \texttt{role: guard} & Quét mã QR ra vào cổng ký túc xá, tra cứu danh sách về muộn, lập biên bản vi phạm giờ giới nghiêm sau 23:00 và trừ điểm kỷ luật KTX. \\ \hline
\textbf{Quản trị viên (Admin)} & \texttt{role: admin} & Quản trị toàn bộ cơ sở dữ liệu, phân quyền tài khoản, cấu hình bảng giá và hệ thống tự động xét duyệt (Auto-Approval Engine). \\ \hline
\end{tabularx}
\end{table}

\section{Danh mục Yêu cầu Chức năng (28 Use Cases)}
Hệ thống bao gồm 28 Use Cases chức năng được tổ chức thành 7 phân hệ chính như trong Bảng~\ref{tab:usecases_summary}.

\begin{table}[htbp]
\centering
\caption{Bảng tổng hợp 28 Use Cases phân theo 7 phân hệ nghiệp vụ}
\label{tab:usecases_summary}
\footnotesize
\begin{tabularx}{\textwidth}{|c|l|l|X|}
\hline
\rowcolor{blue!10}
\textbf{Mã UC} & \textbf{Tên Use Case} & \textbf{Actor chính} & \textbf{Tóm tắt mục tiêu chức năng} \\ \hline
\multicolumn{4}{|l|}{\textbf{Phân hệ 1: Xác thực, Phân quyền \& Hồ sơ cá nhân (Authentication \& Identity)}} \\ \hline
\texttt{UC-01} & Trang chủ giới thiệu KTX & Khách, Sinh viên & Giới thiệu cơ sở vật chất, biểu phí dịch vụ, nội quy chung. \\ \hline
\texttt{UC-02} & Đăng nhập Google OAuth 2.0 & Tất cả & Phân luồng đăng nhập (Gmail sinh viên tuyển sinh vs Email FPT). \\ \hline
\texttt{UC-03} & Quản lý phiên JWT \& Đăng xuất & Tất cả & Quản lý Access/Refresh Token, bảo vệ phiên đăng nhập an toàn. \\ \hline
\texttt{UC-04} & Quản lý Hồ sơ cá nhân (Profile) & Sinh viên & Cập nhật CCCD, thông tin liên hệ khẩn cấp cha mẹ. \\ \hline
\multicolumn{4}{|l|}{\textbf{Phân hệ 2: Quản lý Danh mục \& Cấu trúc Phòng ở (Topology Management)}} \\ \hline
\texttt{UC-05} & Quản lý Tòa nhà \& Tầng lầu & Admin, Manager & Quản lý danh mục tòa nhà (Nam/Nữ) và các tầng lầu. \\ \hline
\texttt{UC-06} & Quản lý Cấu trúc Phòng 1:4:6 & Admin, Manager & Thiết lập 1 phòng = 4 buồng (A,B,C,D), 1 buồng = 6 giường. \\ \hline
\texttt{UC-07} & Cấu hình Bảng giá \& Biểu phí & Admin & Cấu hình giá thuê giường, đơn giá điện, nước, thuế VAT. \\ \hline
\texttt{UC-08} & Quản lý Tài khoản \& Import DSSV & Admin & Quản lý phân quyền RBAC và nạp danh sách sinh viên tuyển sinh. \\ \hline
\multicolumn{4}{|l|}{\textbf{Phân hệ 3: Đặt chỗ Trực quan \& Xét duyệt Tự động (Core Booking Workflow)}} \\ \hline
\texttt{UC-09} & Bản đồ chọn giường trực quan & Sinh viên & Giao diện đồ họa chọn Tòa $\rightarrow$ Tầng $\rightarrow$ Phòng $\rightarrow$ Buồng $\rightarrow$ Giường. \\ \hline
\texttt{UC-10} & Khóa giữ chỗ 15 phút \& VietQR & Sinh viên & Đồng hồ đếm ngược 900 giây và hiển thị mã QR thanh toán cọc. \\ \hline
\texttt{UC-11} & Động cơ Auto-Approval & Hệ thống & Tự động kiểm tra giới tính, công nợ, cọc và duyệt đặt chỗ tức thì. \\ \hline
\texttt{UC-12} & Tạo Hợp đồng điện tử e-Contract & Sinh viên & Tự động render file PDF hợp đồng lưu trú có mã tra cứu duy nhất. \\ \hline
\texttt{UC-13} & Giám sát phòng \& Ngoại lệ đổi chỗ & Manager & Bảng điều khiển quản lý phòng trống và xử lý ngoại lệ đổi giường. \\ \hline
\multicolumn{4}{|l|}{\textbf{Phân hệ 4: Quản lý Công tơ \& Chia đều Tiền Điện Nước (Utility \& Split Billing)}} \\ \hline
\texttt{UC-14} & Quy trình Check-in nhận phòng & Sinh viên, Manager & Lập biên bản bàn giao thiết bị phòng và ghi nhận chỉ số công tơ đầu. \\ \hline
\texttt{UC-15} & Quy trình Check-out trả phòng & Sinh viên, Manager & Nghiệm thu hiện trạng tài sản và lập quyết toán hoàn cọc lưu trú. \\ \hline
\texttt{UC-16} & Ghi số điện nước kèm ảnh chụp & Technician & Form nhập chỉ số $\Delta E, \Delta W$ ngày 25--28 và tải ảnh chụp đồng hồ. \\ \hline
\texttt{UC-17} & Thuật toán chia tiền $Total / N$ & Hệ thống & Tính tổng tiền phòng kèm VAT và chia đều theo đầu sinh viên đang ở. \\ \hline
\texttt{UC-18} & Xem \& Thanh toán Hóa đơn phòng & Sinh viên & Xem hóa đơn minh bạch kèm ảnh công tơ gốc, thanh toán VietQR riêng. \\ \hline
\texttt{UC-19} & Thông báo phát hành \& Nhắc nợ & Hệ thống & Tự động gửi thông báo hóa đơn và nhắc nhở kỳ hạn thanh toán. \\ \hline
\multicolumn{4}{|l|}{\textbf{Phân hệ 5: Tiếp nhận Báo hỏng \& Bảo trì Thiết bị (Maintenance Ticket)}} \\ \hline
\texttt{UC-20} & Nộp phiếu báo hỏng kèm ảnh chụp & Sinh viên & Tạo ticket báo hỏng phân loại Phòng chung vs Buồng cá nhân. \\ \hline
\texttt{UC-21} & Bảng điều phối việc kỹ thuật & Technician, Manager & Nhận phiếu sửa chữa (Bounty Board), phân công công việc. \\ \hline
\texttt{UC-22} & Nghiệm thu \& Ghi nhận vật tư & Technician & Ghi nhận phụ tùng thay thế và tải ảnh kết quả sau sửa chữa. \\ \hline
\texttt{UC-23} & Đánh giá dịch vụ 1--5 sao & Sinh viên & Sinh viên nghiệm thu chất lượng sửa chữa và đánh giá điểm sao. \\ \hline
\multicolumn{4}{|l|}{\textbf{Phân hệ 6: Quản lý Giờ giới nghiêm 23:00 \& Thẻ QR Di động (Curfew \& Security)}} \\ \hline
\texttt{UC-24} & Nộp đơn xin về muộn ngoài giờ & Sinh viên & Tạo đơn xin phép về muộn nộp trước 21:00 kèm lý do hợp lệ. \\ \hline
\texttt{UC-25} & Bảng điều khiển Cổng bảo vệ & Guard & Quét thẻ QR/nhập MSSV, đối chiếu đơn duyệt, phát hiện vi phạm 23:00. \\ \hline
\texttt{UC-26} & Lập biên bản \& Trừ điểm rèn luyện & Guard, Manager & Trừ điểm rèn luyện KTX theo thang 100 điểm, cảnh báo đình chỉ. \\ \hline
\multicolumn{4}{|l|}{\textbf{Phân hệ 7: Báo cáo Thống kê \& Xuất dữ liệu (Analytics \& Export)}} \\ \hline
\texttt{UC-27} & Dashboard tổng quan công suất & Manager & Biểu đồ Donut tỷ lệ lấp đầy giường và biểu đồ cột doanh thu. \\ \hline
\texttt{UC-28} & Xuất báo cáo lưu trú/công nợ & Manager & Xuất danh sách sinh viên tạm trú và nợ tiền sang file Excel (.xlsx). \\ \hline
\end{tabularx}
\end{table}

\section{Đặc tả Chi tiết các Use Case Trọng tâm}

\subsection{Use Case UC-09: Đặt chỗ Giường trực quan (Interactive Bed Booking)}
\begin{itemize}
    \item \textbf{Actor chính:} Sinh viên (Student).
    \item \textbf{Tiền điều kiện:} Sinh viên đã đăng nhập tài khoản thành công và chưa có hợp đồng lưu trú đang hiệu lực trong học kỳ hiện tại.
    \item \textbf{Luồng sự kiện chính (Main Flow):}
    \begin{enumerate}
        \item Sinh viên truy cập trang ``Đăng ký & Chọn Giường'' (\texttt{/booking}).
        \item Hệ thống hiển thị bộ lọc chọn Tòa nhà (mặc định theo giới tính của sinh viên), danh sách Tầng lầu và Phòng trống.
        \item Sinh viên chọn Phòng (ví dụ: \texttt{A205}). Hệ thống hiển thị mô hình không gian của phòng gồm 4 buồng (\texttt{A, B, C, D}).
        \item Sinh viên chọn Buồng mong muốn. Hệ thống hiển thị danh sách 6 giường (\texttt{G01} đến \texttt{G06}) phân loại rõ Tầng dưới (\textit{Lower tier}) và Tầng trên (\textit{Upper tier}) cùng trạng thái màu sắc trực quan (Xanh lá: Còn trống, Xám: Đã có người ở, Đỏ: Đang bảo trì).
        \item Sinh viên click chọn một giường trống (ví dụ: \texttt{A205-B-G03}).
        \item Hệ thống hiển thị bảng tóm tắt chi phí, quy chuẩn phòng và kích hoạt bước Khóa giữ chỗ 15 phút.
    \end{enumerate}
    \item \textbf{Hậu điều kiện:} Giường được chuyển sang trạng thái \texttt{HOLDING}, đồng hồ 15 phút bắt đầu đếm ngược.
\end{itemize}

\subsection{Use Case UC-11: Động cơ Xét duyệt Tự động (Auto-Approval Engine)}
\begin{itemize}
    \item \textbf{Actor chính:} Hệ thống (System Engine).
    \item \textbf{Tiền điều kiện:} Sinh viên gửi yêu cầu xác nhận thanh toán tiền cọc giường thành công.
    \item \textbf{Thuật toán kiểm tra 3 bước:}
    \begin{enumerate}
        \item \textbf{Kiểm tra Giới tính (\textit{Gender Matching}):} Đối chiếu thuộc tính \texttt{user.gender} với \texttt{building.genderType}. Nếu không khớp, từ chối duyệt ngay lập tức để tránh nhầm lẫn phòng nam/nữ.
        \item \textbf{Kiểm tra Lịch sử Tài chính \& Kỷ luật (\textit{Clearance Check}):} Kiểm tra \texttt{user.previousDebt == 0} và \texttt{user.isSuspended == false}. Nếu sinh viên còn nợ tiền điện nước kỳ trước hoặc bị đình chỉ nội trú, yêu cầu giải quyết công nợ trước.
        \item \textbf{Xác thực Tiền cọc (\textit{Deposit Verification}):} Kiểm tra trạng thái giao dịch \texttt{booking.status == 'PAID'} (được đối soát qua cổng VietQR/VNPay).
    \end{enumerate}
    \item \textbf{Kết quả:} Nếu thỏa mãn đủ 3 điều kiện, hệ thống tự động:
    \begin{itemize}
        \item Cập nhật trạng thái giường sang \texttt{OCCUPIED}.
        \item Đổi trạng thái đặt chỗ sang \texttt{APPROVED}.
        \item Khởi tạo bản ghi Hợp đồng lưu trú (\texttt{Contract}) với mã định danh duy nhất và tạo đường dẫn tải file Hợp đồng điện tử PDF.
    \end{itemize}
\end{itemize}

\subsection{Use Case UC-17: Thuật toán Chia đều Tiền Điện Nước ($Total / N$)}
\begin{itemize}
    \item \textbf{Actor chính:} Hệ thống (System Engine).
    \item \textbf{Công thức tính toán:}
    \begin{equation}
        \text{Tiền cá nhân}_i = \frac{T_{\text{Điện}} + T_{\text{Nước}} + T_{\text{Dịch vụ}}}{N_{\text{active}}} \times \omega_i
    \end{equation}
    Trong đó:
    \begin{itemize}
        \item $T_{\text{Điện}} = \Delta E \times \text{Đơn giá điện} \times (1 + \text{VAT}_{\text{điện}})$.
        \item $T_{\text{Nước}} = \Delta W \times \text{Đơn giá nước} \times (1 + \text{VAT}_{\text{nước}})$.
        \item $N_{\text{active}}$: Tổng số sinh viên đang lưu trú thực tế tại phòng trong kỳ hóa đơn.
        \item $\omega_i$: Trọng số tỷ lệ ngày ở (\textit{Pro-rata factor}), với $\omega_i = \frac{\text{Số ngày lưu trú thực tế}}{\text{Tổng số ngày trong tháng}}$. Trường hợp ở trọn vẹn tháng, $\omega_i = 1$.
    \end{itemize}
    \item \textbf{Ràng buộc bảo toàn số tiền:} Sai số làm tròn khi chia cho $N$ sinh viên không được vượt quá $1$ VNĐ so với tổng hóa đơn phòng (sử dụng thuật toán phân bổ phần dư cho sinh viên có số thứ tự đầu tiên).
\end{itemize}

\section{Yêu cầu Phi chức năng (Non-Functional Requirements -- NFR)}
Hệ thống đáp ứng các tiêu chuẩn chất lượng nghiêm ngặt theo ISO/IEC 25010:
\begin{enumerate}
    \item \textbf{Hiệu năng (Performance):} Thời gian phản hồi của các thao tác giao diện không quá $0.5$ giây. Động cơ Auto-Approval hoàn tất quy trình kiểm tra và xuất hợp đồng trong vòng dưới $3$ giây.
    \item \textbf{Khả năng sử dụng (Usability):} Giao diện trực quan, đạt chuẩn Accessibility WCAG 2.1 AA với tỷ lệ tương phản màu sắc cao, hỗ trợ chế độ tương phản đặc biệt (High-contrast mode) cho màn hình di động ngoài trời nắng.
    \item \textbf{Bảo mật (Security):} Sử dụng cơ chế phân quyền RBAC, token JWT kết hợp Refresh Token lưu trong HttpOnly Cookie để chống tấn công XSS/CSRF. Mã QR ra vào cổng áp dụng HMAC mã hóa thay đổi mỗi 60 giây để chống hành vi chụp màn hình chia sẻ cho người khác.
    \item \textbf{Tính thích ứng (Responsive):} Hệ thống tương thích hoàn hảo trên các độ phân giải từ Mobile (375px), Tablet (768px) đến Desktop Full HD (1920px).
\end{enumerate}


\chapter{KIẾN TRÚC HỆ THỐNG VÀ THIẾT KẾ DỮ LIỆU}
\label{ch:architecture}

\section{Kiến trúc Kỹ thuật Tổng thể (System Architecture)}
Hệ thống \textbf{Domanage} được xây dựng theo mô hình kiến trúc Single Page Application (SPA) hiện đại kết hợp với tầng dịch vụ RESTful API được mô phỏng linh hoạt, phân tách rõ ràng giữa tầng hiển thị (Presentation Layer), tầng xử lý nghiệp vụ (Business Logic Layer) và tầng lưu trữ dữ liệu (Data Layer).

\begin{figure}[htbp]
\centering
\begin{tikzpicture}[
    box/.style={rectangle, draw=blue!70!black, fill=blue!5, rounded corners=3pt, thick, align=center, minimum height=1cm},
    layer/.style={rectangle, draw=gray!60, fill=gray!5, dashed, rounded corners=6pt, inner sep=8pt},
    arrow/.style={->, >=stealth, thick, draw=blue!80!black}
]

% Client Layer
\node[box, fill=orange!10, draw=orange!80] (ui) at (0, 4) {\textbf{Presentation Layer (UI/UX)}\\React 19 + TypeScript + Tailwind CSS v4\\Radix UI / ShadCN + Lucide Icons};

% State & Context Layer
\node[box, fill=teal!10, draw=teal!80] (state) at (0, 2.2) {\textbf{State \& Business Logic Layer}\\AuthContext (RBAC) + BookingContext (15m Lock)\\Auto-Approval Engine + Split-Billing Calculator};

% API & Storage Layer
\node[box, fill=blue!10, draw=blue!80] (api) at (0, 0.4) {\textbf{Service \& Mock Data Layer}\\JSON-Server Custom Router (\texttt{server.js})\\RESTful Endpoints + LowDB Local State (\texttt{db.json})};

% External Services
\node[box, fill=purple!10, draw=purple!80] (ext) at (6, 2.2) {\textbf{External Integrations}\\VietQR dynamic Napas247\\jsPDF e-Contract generator\\Excel XML data exporter};

% Arrows
\draw[arrow] (ui) -- (state);
\draw[arrow] (state) -- (api);
\draw[arrow] (state) -- (ext);

\end{tikzpicture}
\caption{Sơ đồ phân tầng kiến trúc phần mềm hệ thống Domanage}
\label{fig:architecture_layers}
\end{figure}

\subsection{Tech Stack chi tiết}
\begin{itemize}
    \item \textbf{Frontend Framework:} React 19 kết hợp cùng TypeScript nghiêm ngặt (Strict Type Safety), mang lại hiệu năng cao và loại trừ lỗi runtime trong quá trình phát triển.
    \item \textbf{Styling \& Component Library:} Tailwind CSS v4 với kiến trúc CSS-first mới nhất, kết hợp các primitives không kiểu dáng từ Radix UI và bộ icon hiện đại Lucide React.
    \item \textbf{Routing \& Phân quyền:} \texttt{react-router-dom} (v7) kết hợp Component bậc cao \texttt{ProtectedRoute} để điều hướng tự động và kiểm soát quyền truy cập theo vai trò.
    \item \textbf{Báo cáo \& Xuất tệp:} \texttt{jspdf} tạo hợp đồng lưu trú điện tử tại máy khách, bộ công cụ tự tạo mã hóa XML xuất báo cáo danh sách tạm trú và công nợ sang định dạng bảng tính Microsoft Excel (\texttt{.xlsx}).
    \item \textbf{Biểu đồ phân tích:} Thư viện biểu đồ trực quan hóa dữ liệu tỷ lệ lấp đầy phòng (Occupancy Donut Chart) và xu hướng doanh thu (Revenue Trend Chart).
\end{itemize}

\section{Mô hình Không gian Phân cấp 1:4:6 (Spatial Topology)}
Một trong những đột phá sáng tạo của dự án là mô hình tổ chức không gian ký túc xá theo tỷ lệ \textbf{1:4:6}, phản ánh chân thực thiết kế phòng tiêu chuẩn tại Đại học FPT:
\begin{itemize}
    \item \textbf{Cấp 1 -- Tòa nhà (Building):} Đại diện cho một khối nhà ký túc xá (ví dụ: Tòa A -- Ký túc xá Nam, Tòa B -- Ký túc xá Nữ).
    \item \textbf{Cấp 2 -- Tầng lầu (Floor):} Mỗi tòa gồm 5 tầng lầu.
    \item \textbf{Cấp 3 -- Phòng ở (Room):} Mỗi tầng bố trí 4 phòng lớn (ví dụ: A101, A102, A103, A104).
    \item \textbf{Cấp 4 -- Buồng sinh hoạt (Compartment):} Mỗi phòng lớn được ngăn thành \textbf{4 buồng cá nhân} riêng biệt, định danh là \textbf{A, B, C, D}.
    \begin{itemize}
        \item \textit{Buồng A (Window Bay):} Hướng ban công/cửa sổ, đón ánh sáng tự nhiên.
        \item \textit{Buồng B (Balcony Wing):} Hướng đón gió mát, gần lối đi nội bộ.
        \item \textit{Buồng C (Study Corner):} Góc học tập yên tĩnh, thoáng đãng.
        \item \textit{Buồng D (Inner Suite):} Không gian ấm cúng, riêng tư tối đa.
    \end{itemize}
    \item \textbf{Cấp 5 -- Giường ngủ (Bed):} Mỗi buồng bao gồm \textbf{6 giường} (tương đương 3 giường tầng: Giường tầng dưới \textit{Lower} và Giường tầng trên \textit{Upper}), được đánh số từ \texttt{G01} đến \texttt{G06}.
\end{itemize}

Như vậy, mỗi phòng lớn có sức chứa tiêu chuẩn là $4 \times 6 = 24$ sinh viên. Toàn bộ mã giường trong hệ thống được chuẩn hóa theo quy tắc chuỗi ký tự bất biến:
\begin{center}
    \texttt{[Mã\_Phòng]-[Mã\_Buồng]-[Mã\_Giường]} $\quad\Longrightarrow\quad$ \textbf{\texttt{A205-B-G03}}
\end{center}

\section{Thiết kế Cơ sở Dữ liệu (Database Schema 12 Thực thể)}
Cơ sở dữ liệu được thiết kế đạt chuẩn chuẩn hóa bậc 3 (3NF) gồm 12 bảng thực thể chính. Chi tiết các thực thể được trình bày trong Bảng~\ref{tab:database_schema}.

\begin{table}[htbp]
\centering
\caption{Bảng từ điển dữ liệu 12 thực thể chính của hệ thống Domanage}
\label{tab:database_schema}
\small
\begin{tabularx}{\textwidth}{|l|l|X|}
\hline
\rowcolor{blue!10}
\textbf{Tên Bảng (Entity)} & \textbf{Khóa chính} & \textbf{Các thuộc tính chính và Mô tả} \\ \hline
\texttt{users} & \texttt{id} & \texttt{mssv}, \texttt{email}, \texttt{name}, \texttt{gender}, \texttt{role}, \texttt{status}, \texttt{previousDebt}, \texttt{isSuspended}, \texttt{disciplinePoints} (thang 100). \\ \hline
\texttt{buildings} & \texttt{id} & \texttt{name}, \texttt{genderType} (MALE/FEMALE), \texttt{totalFloors}. \\ \hline
\texttt{floors} & \texttt{id} & \texttt{buildingId} (FK), \texttt{level} (1--5). \\ \hline
\texttt{rooms} & \texttt{id} & \texttt{floorId} (FK), \texttt{name} (A101--A504), \texttt{status}. \\ \hline
\texttt{compartments} & \texttt{id} & \texttt{roomId} (FK), \texttt{name} (A, B, C, D), \texttt{title}, \texttt{zone}. \\ \hline
\texttt{beds} & \texttt{id} & \texttt{compartmentId} (FK), \texttt{slot} (G01--G06), \texttt{fullCode}, \texttt{tier} (upper/lower), \texttt{status} (available, holding, occupied, maintenance), \texttt{price}. \\ \hline
\texttt{bookings} & \texttt{id} & \texttt{userId} (FK), \texttt{bedId} (FK), \texttt{status} (PENDING, PAID, APPROVED, CANCELLED), \texttt{holdingExpiresAt}. \\ \hline
\texttt{contracts} & \texttt{id} & \texttt{bookingId} (FK), \texttt{userId} (FK), \texttt{bedCode}, \texttt{status} (ACTIVE, TERMINATED), \texttt{validFrom}, \texttt{validTo}. \\ \hline
\texttt{utilities} & \texttt{id} & \texttt{roomId} (FK), \texttt{month}, \texttt{elecOld}, \texttt{elecNew}, \texttt{waterOld}, \texttt{waterNew}, \texttt{elecPhotoUrl}, \texttt{waterPhotoUrl}, \texttt{status}. \\ \hline
\texttt{bills} & \texttt{id} & \texttt{utilityId} (FK), \texttt{userId} (FK), \texttt{amountTotal}, \texttt{status} (UNPAID, PAID), \texttt{qrCodeRef}, \texttt{dueDate}. \\ \hline
\texttt{maintenance\_tickets} & \texttt{id} & \texttt{userId} (FK), \texttt{scope} (ROOM/COMPARTMENT), \texttt{category}, \texttt{urgency}, \texttt{status}, \texttt{technicianId}, \texttt{rating}. \\ \hline
\texttt{gate\_passes} & \texttt{id} & \texttt{userId} (FK), \texttt{curfewDate}, \texttt{reason}, \texttt{expectedReturn}, \texttt{actualReturn}, \texttt{status} (PENDING, APPROVED, VIOLATED). \\ \hline
\end{tabularx}
\end{table}

\section{Vòng đời Trạng thái Giường (Bed State Machine)}
Mỗi vị trí giường ngủ trong hệ thống vận hành theo một máy trạng thái hữu hạn nghiêm ngặt nhằm giải quyết triệt để tình trạng xung đột đặt phòng đồng thời (Concurrency Race Conditions) như minh họa trong Hình~\ref{fig:bed_state_machine}.

\begin{figure}[htbp]
\centering
\begin{tikzpicture}[
    state/.style={circle, draw=blue!80, fill=blue!10, thick, minimum size=1.8cm, align=center, font=\footnotesize\bfseries},
    arrow/.style={->, >=stealth, thick}
]

\node[state, fill=green!15, draw=green!70!black] (avail) at (0, 0) {AVAILABLE\\(Còn trống)};
\node[state, fill=orange!20, draw=orange!80!black] (hold) at (4.5, 0) {HOLDING\\(Khóa 15m)};
\node[state, fill=gray!25, draw=gray!70!black] (occ) at (9, 0) {OCCUPIED\\(Đã ở)};
\node[state, fill=red!15, draw=red!70!black] (maint) at (4.5, -3) {MAINTENANCE\\(Bảo trì)};

\draw[arrow, draw=blue!80!black] (avail) -- node[above, font=\footnotesize] {Click chọn giường} (hold);
\draw[arrow, draw=green!70!black] (hold) -- node[above, font=\footnotesize] {Auto-Approved} (occ);
\draw[arrow, draw=red!70!black, dashed] (hold) to[bend left=30] node[below, font=\footnotesize] {Hết 900s hoặc Hủy} (avail);
\draw[arrow, draw=gray!80!black] (occ) to[bend right=45] node[above, font=\footnotesize] {Check-out trả phòng} (avail);
\draw[arrow, draw=purple!80!black] (avail) -- node[left, font=\footnotesize] {Phát hiện hỏng} (maint);
\draw[arrow, draw=purple!80!black] (maint) -- node[right, font=\footnotesize] {Nghiệm thu sửa xong} (avail);

\end{tikzpicture}
\caption{Sơ đồ chuyển đổi trạng thái của vị trí giường trong hệ thống}
\label{fig:bed_state_machine}
\end{figure}


\chapter{HIỆN THỰC CÁC PHÂN HỆ CỐT LÕI}
\label{ch:core_modules}

\section{Module 1: Xác thực, Phân quyền RBAC và Điều hướng}
Phân hệ xác thực sử dụng React Context API (\texttt{AuthContext}) để quản lý phiên làm việc toàn cục và Component bảo vệ \texttt{ProtectedRoute} để ngăn chặn việc truy cập trái phép vào các đường dẫn nhạy cảm.

\begin{lstlisting}[language=TypeScript, caption={Cơ chế bảo vệ định tuyến theo quyền hạn (ProtectedRoute.tsx)}]
interface ProtectedRouteProps {
  children: React.ReactNode;
  allowedRoles?: Array<"student" | "manager" | "technician" | "admin">;
}

export const ProtectedRoute: React.FC<ProtectedRouteProps> = ({ 
  children, 
  allowedRoles 
}) => {
  const { user, isAuthenticated } = useAuth();
  const location = useLocation();

  if (!isAuthenticated || !user) {
    return <Navigate to="/login" state={{ from: location }} replace />;
  }

  if (allowedRoles && !allowedRoles.includes(user.role)) {
    return <Navigate to="/unauthorized" replace />;
  }

  return <>{children}</>;
};
\end{lstlisting}

Khi người dùng đăng nhập thành công qua giao diện, hệ thống tự động kiểm tra vai trò (\texttt{user.role}) và điều hướng tới phân hệ tương ứng:
\begin{itemize}
    \item Sinh viên (\texttt{student}) $\longrightarrow$ \texttt{/overview} (Bảng điều khiển cá nhân, phòng đang ở, hóa đơn).
    \item Cán bộ Quản lý (\texttt{manager}) $\longrightarrow$ \texttt{/manager} (Dashboard tổng quan, duyệt đơn, xuất báo cáo).
    \item Kỹ thuật viên (\texttt{technician}) $\longrightarrow$ \texttt{/maintenance} (Bảng nhận việc sửa chữa Bounty Board).
\end{itemize}

\section{Module 2 \& 3: Đặt chỗ Trực quan 1:4:6 và Động cơ Auto-Approval}

\subsection{Giao diện Bản đồ Giường trực quan}
Màn hình \texttt{BookingPage.tsx} hiện thực hóa mô hình không gian 1:4:6 cho phép sinh viên duyệt qua từng cấp độ một cách tự nhiên. Khi sinh viên click vào một buồng (ví dụ Buồng B), giao diện hiển thị 6 giường với 4 trạng thái màu sắc trực quan:
\begin{itemize}
    \item \textbf{Màu xanh lá (\texttt{AVAILABLE}):} Giường còn trống, sẵn sàng chọn.
    \item \textbf{Màu cam nhấp nháy (\texttt{HOLDING}):} Giường đang được một sinh viên khác giữ chỗ trong thời gian 15 phút.
    \item \textbf{Màu xám (\texttt{OCCUPIED}):} Giường đã có sinh viên đang sinh sống.
    \item \textbf{Màu đỏ (\texttt{MAINTENANCE}):} Giường đang hỏng hóc hoặc khóa bảo trì.
\end{itemize}

\subsection{Cơ chế Khóa giữ chỗ 15 phút (Countdown Timer)}
Khi sinh viên chọn giường trống, hệ thống kích hoạt bộ đếm thời gian 900 giây (15 phút). Nếu quá thời gian đếm ngược mà sinh viên chưa xác nhận thanh toán cọc, sự kiện \texttt{rollback} sẽ tự động kích hoạt để giải phóng giường về trạng thái \texttt{AVAILABLE}.

\begin{lstlisting}[language=TypeScript, caption={Đồng hồ đếm ngược giữ chỗ 15 phút và tự động giải phóng}]
useEffect(() => {
  if (!isLocked) return;
  const timer = setInterval(() => {
    setTimeLeft((prev) => {
      if (prev <= 1) {
        clearInterval(timer);
        setIsLocked(false);
        toast.warning("Hết thời gian giữ chỗ! Giường đã được giải phóng.");
        return 0;
      }
      return prev - 1;
    });
  }, 1000);
  return () => clearInterval(timer);
}, [isLocked]);
\end{lstlisting}

\subsection{Động cơ Tự động Xét duyệt (Auto-Approval Engine)}
Tại tầng backend mô phỏng (\texttt{server.js}), API \texttt{/api/bookings/approve} tiếp nhận mã đặt phòng và tự động kiểm tra 3 điều kiện tiên quyết mà không cần bất kỳ sự can thiệp thủ công nào từ phía con người:

\begin{lstlisting}[language=JavaScript, caption={Thuật toán 3 bước xét duyệt tự động trong server.js}]
server.post('/api/bookings/approve', (req, res) => {
  const { bookingId } = req.body;
  const db = router.db;

  const booking = db.get('bookings').find({ id: bookingId }).value();
  const user = db.get('users').find({ id: booking.userId }).value();
  const bed = db.get('beds').find({ id: booking.bedId }).value();
  const building = getBuildingByBedId(db, bed.id);

  // Dieu kien 1: Kiem tra khop gioi tinh
  if (user.gender !== building.genderType) {
    return res.status(400).json({ success: false, message: 'Sai gioi tinh toa nha.' });
  }

  // Dieu kien 2: Kiem tra cong no cu & ky luat
  if (user.previousDebt > 0 || user.isSuspended) {
    return res.status(400).json({ success: false, message: 'Sinh vien con no hoac bi dinh chi.' });
  }

  // Dieu kien 3: Xac thuc thanh toan 100% tien coc
  if (booking.status !== 'PAID') {
    return res.status(400).json({ success: false, message: 'Chua thanh toan tien coc.' });
  }

  // Phe duyet tu dong & Khoi tao Hop dong dien tu
  db.get('bookings').find({ id: bookingId }).assign({ status: 'APPROVED' }).write();
  db.get('beds').find({ id: booking.bedId }).assign({ status: 'OCCUPIED', studentId: user.id }).write();

  const contract = {
    id: `ct_${Date.now()}`,
    bookingId,
    userId: user.id,
    bedCode: bed.fullCode,
    status: 'ACTIVE',
    createdAt: new Date().toISOString()
  };
  db.get('contracts').push(contract).write();

  return res.json({ success: true, message: 'Auto-Approval hoan tat', data: { contract } });
});
\end{lstlisting}

Sau khi API duyệt thành công, giao diện phía client tự động kích hoạt thư viện \texttt{jspdf} để xuất file PDF Hợp đồng lưu trú gửi đến thiết bị của sinh viên với đầy đủ thông tin: họ tên, MSSV, mã giường, giá thuê, điều khoản ký túc xá và mã xác thực điện tử duy nhất.

\section{Module 4: Quản lý Công tơ và Thuật toán Chia đều Tiền Điện Nước}

\subsection{Giao diện Ghi chỉ số Công tơ định kỳ}
Vào ngày 25--28 hàng tháng, Kỹ thuật viên tòa nhà sử dụng phân hệ \texttt{BillingPage.tsx} để nhập chỉ số điện mới ($\text{elecNew}$) và chỉ số nước mới ($\text{waterNew}$). Điểm đặc biệt là hệ thống bắt buộc kỹ thuật viên phải tải lên hình ảnh chụp thực tế đồng hồ điện và đồng hồ nước tại cửa phòng. Hình ảnh này được lưu trữ và hiển thị trực tiếp trên hóa đơn của sinh viên, chấm dứt hoàn toàn tình trạng nghi ngờ gian lận hoặc sai lệch chỉ số.

\subsection{Thuật toán Phân bổ Hóa đơn ($Total / N$)}
Hệ thống tính toán tổng tiền phòng và phân bổ cho từng sinh viên:
\begin{lstlisting}[language=TypeScript, caption={Thuật toán tính toán và chia đều hóa đơn phòng}]
const calculateSplitBill = (
  elecUsage: number,     // Delta E (kWh)
  waterUsage: number,    // Delta W (m3)
  elecRate: number = 3200,
  waterRate: number = 18000,
  vatRate: number = 0.10,
  cleaningFee: number = 120000,
  residentsCount: number = 6
) => {
  const elecCost = elecUsage * elecRate * (1 + vatRate);
  const waterCost = waterUsage * waterRate * (1 + vatRate);
  const totalRoomBill = elecCost + waterCost + cleaningFee;

  // Chia deu cho N sinh vien dang luu tru
  const perStudentBase = Math.floor(totalRoomBill / residentsCount);
  const remainder = totalRoomBill - (perStudentBase * residentsCount);

  return {
    elecCost,
    waterCost,
    totalRoomBill,
    perStudentAmount: perStudentBase,
    remainderCorrection: remainder // Cong vao sinh vien stt 1 de bao toan tong
  };
};
\end{lstlisting}

Mỗi sinh viên khi đăng nhập vào hệ thống sẽ thấy hóa đơn cá nhân với mã VietQR riêng biệt chứa nội dung chuyển khoản được định dạng chuẩn: \texttt{DMS [Mã\_HĐ] [MSSV]}, giúp việc gạch nợ diễn ra hoàn toàn tự động.

\section{Module 5: Phân hệ Tiếp nhận Báo hỏng và Bảo trì Thiết bị}
Quy trình bảo trì thiết bị ký túc xá được khép kín thành vòng đời 5 trạng thái thông qua \texttt{MaintenancePage.tsx}:
\begin{enumerate}
    \item \textbf{Khởi tạo (SUBMITTED):} Sinh viên nộp phiếu báo hỏng, chọn phạm vi (Thiết bị phòng chung như điều hòa, bình nóng lạnh hoặc Thiết bị buồng riêng như bóng đèn bàn học, ổ cắm), chụp ảnh hiện trạng hỏng hóc và chọn mức độ khẩn cấp (Thường, Gấp, Khẩn cấp).
    \item \textbf{Tiếp nhận (ASSIGNED):} Phiếu được đưa lên \textit{Technician Bounty Board}. Kỹ thuật viên tòa nhà chủ động nhận việc hoặc được Ban Quản lý điều phối.
    \item \textbf{Đang xử lý (IN\_PROGRESS):} Kỹ thuật viên đến hiện trường kiểm tra, thay thế phụ tùng và cập nhật ghi chú kỹ thuật.
    \item \textbf{Nghiệm thu hoàn thành (COMPLETED):} Kỹ thuật viên chụp ảnh hiện trạng sau khi sửa chữa xong, nhập danh mục linh kiện thay thế và đánh dấu hoàn thành.
    \item \textbf{Đánh giá dịch vụ (REVIEWED):} Sinh viên nhận được thông báo, xác nhận thiết bị hoạt động tốt và chấm điểm mức độ hài lòng từ 1 đến 5 sao.
\end{enumerate}

\section{Module 6: Thẻ QR Di động và Kiểm soát Giờ Giới Nghiêm 23:00}
Phân hệ \texttt{MobileAppPage.tsx} đóng vai trò như một ứng dụng di động giả lập dành cho sinh viên và bảo vệ trực cổng:
\begin{itemize}
    \item \textbf{Thẻ sinh viên điện tử (Digital Gate Pass):} Hiển thị mã QR động được sinh dựa trên thuật toán mã hóa HMAC. Mã QR này tự động làm mới sau mỗi 60 giây và hiển thị đồng hồ đếm ngược, ngăn chặn việc chụp ảnh màn hình gửi cho người khác đi hộ.
    \item \textbf{Đơn xin về muộn ngoài giờ:} Sinh viên có lịch học thêm hoặc làm thêm phải nộp đơn xin phép trước 21:00. Đơn được Quản lý KTX duyệt điện tử.
    \item \textbf{Bảng điều khiển Bảo vệ (Security Gate Console):} Khi sinh viên qua cổng sau 23:00, bảo vệ quét mã QR hoặc nhập nhanh MSSV. Hệ thống lập tức đối chiếu:
    \begin{itemize}
        \item Nếu sinh viên \textbf{có đơn xin về muộn đã duyệt}: Cổng mở, hệ thống ghi nhận trạng thái về hợp lệ.
        \item Nếu sinh viên \textbf{không có đơn}: Hệ thống bật cảnh báo đỏ, tự động ghi nhận biên bản vi phạm giờ giới nghiêm và trừ 15 điểm rèn luyện KTX. Nếu tổng điểm rèn luyện giảm dưới 70 điểm, sinh viên bị cảnh cáo; dưới 50 điểm sẽ bị đình chỉ quyền nội trú học kỳ kế tiếp.
    \end{itemize}
\end{itemize}

\section{Module 7: Bảng điều khiển Quản trị và Báo cáo Excel}
Màn hình \texttt{ManagerDashboardPage.tsx} cung cấp các chỉ số kinh doanh quan trọng (KPIs) cho Ban Quản lý KTX:
\begin{itemize}
    \item \textbf{Biểu đồ Donut Tỷ lệ Lấp đầy (Occupancy Rate):} Thống kê số giường đang ở, số giường trống và số giường bảo trì theo từng tầng, từng tòa nhà.
    \item \textbf{Biểu đồ Cột Xu hướng Doanh thu (Revenue Trend):} Theo dõi doanh thu tiền phòng và tiền dịch vụ điện nước theo từng tháng.
    \item \textbf{Xuất báo cáo định dạng Excel (\texttt{.xlsx}):} Tích hợp module \texttt{excelExport.ts} cho phép xuất danh sách sinh viên tạm trú (phục vụ đăng ký lưu trú với Công an địa phương) và danh sách sinh viên nợ tiền quá hạn chỉ với một cú click chuột.
\end{itemize}


\chapter{ĐẢM BẢO CHẤT LƯỢNG VÀ KIỂM THỬ PHẦN MỀM}
\label{ch:testing}

\section{Chiến lược Kiểm thử và Phương pháp luận (QA Strategy)}
Công tác đảm bảo chất lượng phần mềm (QA) được thực hiện xuyên suốt các giai đoạn phát triển dự án bởi Kỹ sư kiểm thử Huỳnh Gia Kim, áp dụng mô hình kiểm thử hộp đen (Black-box Testing) kết hợp kiểm thử tích hợp luồng nghiệp vụ người dùng (End-to-End User Journeys).

Hệ thống được kiểm thử dựa trên 3 tiêu chí cốt lõi:
\begin{enumerate}
    \item \textbf{Kiểm thử chức năng (Functional Testing):} Đảm bảo toàn bộ 28 Use Cases hoạt động chính xác theo yêu cầu nghiệp vụ đã đặc tả trong tài liệu SRS.
    \item \textbf{Kiểm thử biên và ngoại lệ (Edge-case \& Boundary Testing):} Đưa ra các tình huống tải cao, xung đột đồng thời và dữ liệu bất thường để kiểm tra tính bền bỉ của hệ thống.
    \item \textbf{Kiểm thử giao diện \& tương thích (UI/UX \& Cross-device Testing):} Kiểm tra khả năng tương thích trên các kích thước màn hình phổ biến (Mobile, Tablet, Desktop) và các trình duyệt web hiện đại.
\end{enumerate}

\section{Kiểm thử 4 Ca Biên Then chốt (Critical Edge-cases)}

\subsection{Case 1: Xung đột Đặt phòng Đồng thời (Concurrency Race Condition)}
\begin{itemize}
    \item \textbf{Mô tả kịch bản:} Hai sinh viên (Sinh viên A và Sinh viên B) cùng truy cập buồng B205 cùng một thời điểm và đồng thời click chọn vị trí giường trống \texttt{A205-B-G03} tại cùng một giây.
    \item \textbf{Hành vi kỳ vọng:} Hệ thống chỉ cho phép đúng 1 sinh viên (người gửi tín hiệu đến trước) khóa giường thành công (\texttt{HOLDING}); sinh viên còn lại nhận được thông báo: \textit{``Rất tiếc! Giường này vừa được giữ chỗ bởi sinh viên khác, vui lòng chọn giường khác''}, và vị trí giường lập tức chuyển sang màu cam trên màn hình của sinh viên thứ hai.
    \item \textbf{Kết quả kiểm thử:} \textbf{\color{green!70!black}ĐẠT (PASS)}. Cơ chế khóa tức thời và cập nhật trạng thái đã giải quyết trơn tru xung đột.
\end{itemize}

\subsection{Case 2: Hết hạn Giữ chỗ 15 phút (Timeout Rollback)}
\begin{itemize}
    \item \textbf{Mô tả kịch bản:} Sinh viên chọn giường và bắt đầu đếm ngược 15 phút (900 giây) nhưng đóng trình duyệt hoặc không thực hiện quét mã VietQR để thanh toán cọc.
    \item \textbf{Hành vi kỳ vọng:} Khi đồng hồ đếm về 0, hệ thống kích hoạt cơ chế Rollback, tự động giải phóng bản ghi giữ chỗ (\texttt{CANCELLED}), đưa trạng thái giường từ \texttt{HOLDING} trở lại \texttt{AVAILABLE} (màu xanh lá) để các sinh viên khác có thể đăng ký bình thường.
    \item \textbf{Kết quả kiểm thử:} \textbf{\color{green!70!black}ĐẠT (PASS)}. Giường được trả về trạng thái trống chính xác ở giây thứ 901.
\end{itemize}

\subsection{Case 3: Sai số Làm tròn khi Chia tiền ($Total / N$)}
\begin{itemize}
    \item \textbf{Mô tả kịch bản:} Tổng tiền điện nước của phòng là $1.000.000$ VNĐ, phòng có 6 sinh viên đang ở ($1.000.000 / 6 = 166.666,666...$ VNĐ). Nếu làm tròn xuống $166.666$ VNĐ thì tổng thu chỉ là $999.996$ VNĐ (thiếu $4$ VNĐ). Nếu làm tròn lên $166.667$ VNĐ thì tổng thu là $1.000.002$ VNĐ (thừa $2$ VNĐ).
    \item \textbf{Hành vi kỳ vọng:} Thuật toán \texttt{calculateSplitBill} thực hiện phân bổ phần dư (\textit{Remainder Correction}): 4 sinh viên đầu tiên đóng $166.667$ VNĐ, 2 sinh viên còn lại đóng $166.666$ VNĐ. Đảm bảo tổng tiền thu từ các cá nhân khớp chính xác $100\%$ ($1.000.000$ VNĐ), không lệch dù chỉ $1$ đồng.
    \item \textbf{Kết quả kiểm thử:} \textbf{\color{green!70!black}ĐẠT (PASS)}. Hoàn toàn triệt tiêu sai số làm tròn số học.
\end{itemize}

\subsection{Case 4: Quét cổng sau 23:00 không có đơn xin phép}
\begin{itemize}
    \item \textbf{Mô tả kịch bản:} Sinh viên quét mã QR tại cổng ký túc xá lúc 23:15 nhưng không nộp đơn xin về muộn trước 21:00.
    \item \textbf{Hành vi kỳ vọng:} Giao diện bảo vệ lập tức hiển thị cảnh báo vi phạm màu đỏ, âm thanh cảnh báo phát ra, hệ thống tự động ghi biên bản vi phạm giờ giới nghiêm kèm mốc thời gian thực và tự động trừ 15 điểm rèn luyện của sinh viên.
    \item \textbf{Kết quả kiểm thử:} \textbf{\color{green!70!black}ĐẠT (PASS)}. Hệ thống ghi nhận biên bản vi phạm chính xác và cập nhật điểm rèn luyện ngay lập tức.
\end{itemize}

\section{Bảng Tổng hợp Kết quả Kiểm thử Toàn bộ Phân hệ}
Bảng~\ref{tab:test_execution_matrix} tóm tắt kết quả kiểm thử thực nghiệm trên toàn bộ 7 phân hệ của hệ thống Domanage.

\begin{table}[htbp]
\centering
\caption{Ma trận kết quả thực thi kiểm thử 28 Use Cases}
\label{tab:test_execution_matrix}
\small
\begin{tabularx}{\textwidth}{|l|c|c|c|c|}
\hline
\rowcolor{blue!10}
\textbf{Phân hệ Chức năng (Module)} & \textbf{Tổng số Cases} & \textbf{Passed} & \textbf{Failed} & \textbf{Tỷ lệ Đạt} \\ \hline
Module 1: Xác thực, Phân quyền \& Profile & 12 & 12 & 0 & 100\% \\ \hline
Module 2: Quản lý Danh mục \& Cấu trúc 1:4:6 & 8 & 8 & 0 & 100\% \\ \hline
Module 3: Đặt chỗ Trực quan \& Auto-Approval & 18 & 18 & 0 & 100\% \\ \hline
Module 4: Công tơ \& Chia tiền Điện Nước & 16 & 16 & 0 & 100\% \\ \hline
Module 5: Phiếu Báo hỏng \& Bảo trì Thiết bị & 14 & 14 & 0 & 100\% \\ \hline
Module 6: Giờ Giới Nghiêm 23:00 \& Thẻ QR Cổng & 12 & 12 & 0 & 100\% \\ \hline
Module 7: Dashboard Quản trị \& Xuất Excel & 10 & 10 & 0 & 100\% \\ \hline
\rowcolor{green!10}
\textbf{TỔNG CỘNG} & \textbf{90} & \textbf{90} & \textbf{0} & \textbf{100\%} \\ \hline
\end{tabularx}
\end{table}

\section{Kiểm thử Đa Trình duyệt và Tính tương thích Thiết bị}
Hệ thống đã trải qua quy trình kiểm thử khả năng hiển thị tương thích (Responsive Testing) trên các môi trường thực nghiệm đa dạng:
\begin{itemize}
    \item \textbf{Màn hình Mobile (iPhone 13, 14, Samsung Galaxy S22 -- độ phân giải 375px -- 414px):} Toàn bộ thanh điều hướng được chuyển thành thanh điều hướng đáy (\textit{Bottom Mobile Navigation Bar}), bản đồ giường tự động co giãn theo trục dọc (Vertical scroll), thao tác chạm vuốt mượt mà.
    \item \textbf{Màn hình Tablet (iPad Pro, iPad Air -- độ phân giải 768px -- 1024px):} Bố cục hiển thị lưới 2 cột, bảng điều khiển bảo vệ và kỹ thuật viên hiển thị rõ ràng, dễ thao tác.
    \item \textbf{Màn hình Desktop (Full HD 1920x1080):} Tận dụng tối đa không gian màn hình với Sidebar cố định, bản đồ buồng 1:4:6 hiển thị toàn cảnh 4 buồng A, B, C, D trên một khung nhìn duy nhất.
    \item \textbf{Trình duyệt hỗ trợ:} Đã kiểm thử và chạy ổn định $100\%$ trên Google Chrome (v128+), Microsoft Edge (v128+), Mozilla Firefox (v130+) và Apple Safari (v17+).
\end{itemize}


\chapter{QUẢN TRỊ DỰ ÁN VÀ PHÂN CÔNG CÔNG VIỆC}
\label{ch:project_management}

\section{Cơ cấu Tổ chức Nhóm và Ma trận Trách nhiệm RACI}
Dự án được thực hiện bởi nhóm gồm 6 thành viên với sự phân công vai trò rõ ràng, chuyên môn hóa cao theo từng phân hệ nghiệp vụ. Để quản trị trách nhiệm hiệu quả, nhóm áp dụng mô hình \textbf{Ma trận RACI} (Responsible, Accountable, Consulted, Informed) như được trình bày chi tiết trong Bảng~\ref{tab:raci_matrix}.

\begin{table}[htbp]
\centering
\caption{Ma trận Trách nhiệm RACI chi tiết 6 thành viên nhóm FER202\_Group4}
\label{tab:raci_matrix}
\small
\begin{tabularx}{\textwidth}{|l|c|c|c|c|c|c|}
\hline
\rowcolor{blue!10}
\textbf{Phân hệ / Gói công việc (WBS)} & \textbf{Vinh} & \textbf{Tiến} & \textbf{Nguyên} & \textbf{V.Trung} & \textbf{Đ.Trung} & \textbf{Kim} \\ \hline
Module 1: Authentication \& Identity (OAuth, JWT) & \textbf{R/A} & C & C & I & I & C \\ \hline
Module 2: Topology Cấu trúc Không gian 1:4:6 & A & \textbf{R} & \textbf{R} & I & I & C \\ \hline
Module 3: Đặt giường Trực quan \& Auto-Approval & A & \textbf{R} & \textbf{R} & I & I & C \\ \hline
Module 4: Chia tiền Điện Nước ($Total / N$) \& Handover & A & C & C & \textbf{R} & I & C \\ \hline
Module 5: Phiếu Báo hỏng \& Bảo trì Thiết bị & A & I & I & I & \textbf{R} & C \\ \hline
Module 6: Giờ Giới Nghiêm 23:00 \& Thẻ QR Cổng & A & I & I & I & \textbf{R} & C \\ \hline
Module 7: Dashboard Quản trị \& Xuất Báo cáo Excel & A & C & I & I & \textbf{R} & C \\ \hline
Đảm bảo Chất lượng, Viết Test Cases \& Soạn Slide & A & I & I & I & I & \textbf{R} \\ \hline
\end{tabularx}
\end{table}

\begin{flushleft}
\footnotesize
\textit{Quy ước:} \textbf{R (Responsible):} Người trực tiếp thực hiện; \textbf{A (Accountable):} Người chịu trách nhiệm cao nhất (Lead); \textbf{C (Consulted):} Người tham gia tư vấn/hỗ trợ; \textbf{I (Informed):} Người nhận thông tin kết quả.
\end{flushleft}

\section{Lộ trình Triển khai 5 Tuần (Sprint Roadmap)}
Quá trình phát triển dự án tuân thủ phương pháp luận linh hoạt (Agile Scrum), chia thành 5 Sprint tuần rõ ràng:
\begin{itemize}
    \item \textbf{Tuần 1 -- Thiết lập Kiến trúc \& Nền tảng (Foundation Week):}
    \begin{itemize}
        \item Khảo sát nghiệp vụ ký túc xá, hoàn thiện tài liệu SRS và đặc tả 28 Use Cases.
        \item Cấu hình mã nguồn Git Repository, tích hợp Vite, TypeScript, Tailwind CSS v4, Prettier, ESLint.
        \item Thiết kế kiến trúc định tuyến và State Management (\texttt{AuthContext}, \texttt{BookingContext}).
        \item Xây dựng cơ sở dữ liệu mô phỏng 12 thực thể (\texttt{db.json}) và nạp dữ liệu mẫu Tòa A (5 tầng, 20 phòng).
    \end{itemize}

    \item \textbf{Tuần 2 -- Luồng Đặt phòng Trọng tâm (Core Booking Workflow):}
    \begin{itemize}
        \item Xây dựng Component bản đồ phòng trực quan 1:4:6 (Tòa $\rightarrow$ Tầng $\rightarrow$ Phòng $\rightarrow$ Buồng $\rightarrow$ Giường).
        \item Tích hợp Modal khóa giữ chỗ 15 phút, đồng hồ đếm ngược và hiển thị cổng VietQR động.
        \item Cài đặt động cơ tự động xét duyệt (Auto-Approval Engine) 3 bước trong \texttt{server.js}.
        \item Tích hợp thư viện \texttt{jspdf} tự động kết xuất Hợp đồng lưu trú điện tử PDF.
    \end{itemize}

    \item \textbf{Tuần 3 -- Điện Nước \& Quy trình Bàn giao (Billing \& Handover Workflow):}
    \begin{itemize}
        \item Xây dựng form ghi chỉ số công tơ điện nước định kỳ cho Kỹ thuật viên (kèm tính năng tải ảnh đồng hồ).
        \item Cài đặt thuật toán tự động chia tiền phòng $Total / N$ kèm xử lý sai số làm tròn số học và tính tiền pro-rata.
        \item Xây dựng giao diện Hóa đơn cá nhân minh bạch kèm mã VietQR cá thể hóa.
        \item Hiện thực form biên bản điện tử cho quy trình Check-in và Check-out bàn giao tài sản.
    \end{itemize}

    \item \textbf{Tuần 4 -- Vận hành Cơ sở vật chất \& Quản trị (Operations \& Analytics):}
    \begin{itemize}
        \item Hoàn thiện hệ thống tiếp nhận phiếu báo hỏng và bảng nhận việc kỹ thuật viên (Bounty Board).
        \item Xây dựng phân hệ Thẻ sinh viên điện tử QR động (HMAC làm mới mỗi 60s) và Bảng điều khiển Cổng bảo vệ.
        \item Cài đặt thuật toán phát hiện vi phạm giờ giới nghiêm 23:00 và tự động trừ điểm rèn luyện KTX.
        \item Xây dựng Dashboard quản trị với biểu đồ tỷ lệ lấp đầy phòng và tính năng xuất dữ liệu tạm trú/công nợ sang Excel.
    \end{itemize}

    \item \textbf{Tuần 5 -- Kiểm thử Toàn diện \& Chuẩn bị Bảo vệ (Final Polish \& Defense):}
    \begin{itemize}
        \item Thực thi 90 ca kiểm thử cho toàn bộ 28 Use Cases và kiểm tra các kịch bản biên (Edge-cases).
        \item Tối ưu hóa hiệu năng tải trang, kiểm thử Responsive đa thiết bị và sửa chữa toàn bộ lỗi tồn đọng.
        \item Đóng băng mã nguồn (Code Freeze), viết tài liệu báo cáo dự án và chuẩn bị kịch bản demo bảo vệ.
    \end{itemize}
\end{itemize}

\section{Quy ước Làm việc Nhóm (Team Working Agreements)}
Để đảm bảo tính đồng bộ và chất lượng mã nguồn cao nhất, nhóm đã thống nhất và tuân thủ các quy ước kỹ thuật sau:
\begin{enumerate}
    \item \textbf{Chiến lược Nhánh Git (Git Flow):}
    \begin{itemize}
        \item Nhánh \texttt{main}: Mã nguồn phiên bản phát hành chính thức, chỉ merge khi có sự chấp thuận của Project Manager (Vinh).
        \item Nhánh \texttt{develop}: Nhánh tích hợp liên tục của cả nhóm.
        \item Các nhánh tính năng con: Đặt tên theo tiền tố chuẩn: \texttt{feature/booking-146}, \texttt{feature/utility-split}, \texttt{feature/curfew-guard}, v.v.
    \end{itemize}
    \item \textbf{Quy chuẩn Thông điệp Commit (Conventional Commits):}
    \begin{itemize}
        \item \texttt{feat: thêm chức năng chọn giường trực quan 1:4:6}
        \item \texttt{fix: khắc phục lỗi đếm ngược 15 phút không nhả giường}
        \item \texttt{style: tinh chỉnh giao diện hóa đơn cá nhân trên mobile}
        \item \texttt{docs: cập nhật tài liệu kiểm thử và báo cáo LaTeX}
    \end{itemize}
    \item \textbf{Giao tiếp và Đánh giá:} Nhóm duy trì phiên họp nhanh (Daily Standup) 15 phút vào mỗi sáng để rà soát tiến độ và giải quyết ngay các điểm nghẽn kỹ thuật.
\end{enumerate}

\section{Đánh giá Mức độ Đóng góp của Từng Thành viên}
Tất cả 6 thành viên trong nhóm đã hoàn thành xuất sắc các nhiệm vụ được giao, đóng góp tích cực vào sự thành công chung của đề tài như được ghi nhận trong Bảng~\ref{tab:member_evaluation}.

\begin{table}[htbp]
\centering
\caption{Bảng đánh giá mức độ hoàn thành nhiệm vụ của các thành viên}
\label{tab:member_evaluation}
\small
\begin{tabularx}{\textwidth}{|c|l|X|c|c|}
\hline
\rowcolor{blue!10}
\textbf{STT} & \textbf{Họ và Tên} & \textbf{Sản phẩm bàn giao cụ thể (Deliverables)} & \textbf{Tiến độ} & \textbf{Đóng góp} \\ \hline
1 & Nguyễn Hoàng Vinh & Kiến trúc dự án, AuthContext, ProtectedRoute, tích hợp API, quản lý tiến độ. & 100\% & 100\% \\ \hline
2 & Trần Công Tiến & Design System, Bản đồ chọn giường 1:4:6, Modal giữ chỗ 15m, VietQR UI. & 100\% & 100\% \\ \hline
3 & Trần Thanh Nguyên & Schema 12 bảng, Seed Data, Auto-Approval Engine, kết xuất Hợp đồng PDF. & 100\% & 100\% \\ \hline
4 & Văn Thái Trung & Form ghi chỉ số công tơ kèm ảnh, Thuật toán chia tiền $Total/N$, Form Handover. & 100\% & 100\% \\ \hline
5 & Đinh Thanh Trung & Ticket báo hỏng, Bounty Board, Quản lý giới nghiêm 23:00, Dashboard Recharts, Xuất Excel. & 100\% & 100\% \\ \hline
6 & Huỳnh Gia Kim & Bộ 90 Test Cases, kiểm thử 4 Edge-cases, kiểm thử Responsive, Slide \& Kịch bản demo. & 100\% & 100\% \\ \hline
\end{tabularx}
\end{table}


\chapter{KẾT LUẬN VÀ HƯỚNG PHÁT TRIỂN}
\label{ch:conclusion}

\section{Kết quả Đạt được của Đề tài}
Trải qua 5 tuần nghiên cứu và phát triển nghiêm túc, nhóm sinh viên \textbf{FER202\_Group4} đã hoàn thành toàn diện các mục tiêu đề ra cho dự án \textbf{Domanage -- Hệ thống Quản lý Ký túc xá Trực tuyến}:
\begin{enumerate}
    \item \textbf{Về mặt Học thuật và Kỹ thuật:}
    \begin{itemize}
        \item Làm chủ và ứng dụng thành công các công nghệ tiên tiến nhất trong hệ sinh thái Front-end hiện đại: React 19, TypeScript, Tailwind CSS v4, Radix UI và Vite.
        \item Xây dựng cấu trúc mã nguồn dạng module phân tách cao, tuân thủ nguyên lý Clean Architecture và Single Responsibility Principle (SRP).
        \item Thiết kế và hiện thực hóa thành công \textbf{Mô hình không gian 1:4:6} và \textbf{Động cơ Tự động xét duyệt (Auto-Approval Engine)}, giải quyết triệt để bài toán xếp phòng và phê duyệt hợp đồng lưu trú vốn tốn nhiều thời gian và nhân lực.
    \end{itemize}

    \item \textbf{Về mặt Nghiệp vụ và Thực tiễn:}
    \begin{itemize}
        \item Hệ thống giải quyết trọn vẹn 28 Use Cases nghiệp vụ trải dài qua 7 phân hệ chức năng: từ đăng ký phòng, hợp đồng điện tử, quản lý công tơ, tính toán chia tiền điện nước minh bạch, bảo trì tài sản cho tới giám sát an ninh giờ giới nghiêm.
        \item Xây dựng giao diện Responsive tối ưu, cho phép cán bộ quản lý làm việc hiệu quả trên máy tính bàn và sinh viên, nhân viên bảo vệ thao tác mượt mà trên thiết bị di động.
    \end{itemize}

    \item \textbf{Về mặt Đảm bảo Chất lượng và Quản trị Dự án:}
    \begin{itemize}
        \item Hoàn thành bộ 90 ca kiểm thử thực nghiệm đạt tỷ lệ thành công 100\%, xử lý triệt để các tình huống xung đột đồng thời, sai số làm tròn số học và khôi phục trạng thái khi hết hạn thời gian giữ chỗ.
        \item Đội ngũ 6 thành viên phối hợp nhịp nhàng theo quy trình Agile Scrum, bảo đảm tiến độ bàn giao sản phẩm chuẩn xác theo từng mốc Sprint.
    \end{itemize}
\end{enumerate}

\section{Hạn chế của Hệ thống Hiện tại}
Bên cạnh những kết quả nổi bật đã đạt được, dự án vẫn còn một số hạn chế do giới hạn về thời gian và phạm vi của môn học Front-End Web:
\begin{itemize}
    \item \textbf{Tầng Backend:} Hệ thống hiện tại đang sử dụng máy chủ mô phỏng JSON-Server kết hợp LowDB; chưa triển khai một hệ thống Backend hoàn chỉnh bằng Node.js (NestJS/Express) hoặc Spring Boot kết nối PostgreSQL phân tán trong môi trường sản xuất thực tế.
    \item \textbf{Cổng thanh toán:} Các giao dịch thanh toán VietQR và VNPay hiện tại đang dừng ở mức độ mô phỏng giao thức và tạo mã QR động; chưa kết nối Webhook trực tiếp với Ngân hàng thương mại để tự động lắng nghe biến động số dư theo thời gian thực.
    \item \textbf{Hạ tầng phần cứng:} Việc ghi nhận chỉ số công tơ điện nước vẫn cần kỹ thuật viên đến tận nơi đọc và chụp ảnh thủ công, chưa được tự động hóa bằng cảm biến IoT.
\end{itemize}

\section{Hướng Phát triển trong Tương lai}
Nhằm phát triển Domanage thành một giải pháp phần mềm thương mại hoàn chỉnh cho các trường đại học và khu đô thị đại học, nhóm đề xuất các hướng nâng cấp tiếp theo:
\begin{enumerate}
    \item \textbf{Phát triển Ứng dụng Di động Native (React Native / Flutter):} Chuyển đổi phân hệ Mobile Web sang ứng dụng di động độc lập trên App Store và Google Play, tích hợp Push Notification đẩy thông báo hóa đơn và nhắc nhở giờ giới nghiêm theo thời gian thực.
    \item \textbf{Tích hợp Hệ sinh thái Vạn vật Kết nối (IoT Smart Dormitory):}
    \begin{itemize}
        \item Lắp đặt công tơ điện và đồng hồ nước thông minh có kết nối Wi-Fi/LoRaWAN để tự động truyền dữ liệu tiêu thụ về hệ thống mỗi ngày, phát hiện cảnh báo rò rỉ nước hoặc chập cháy điện tức thời.
        \item Lắp đặt khóa cửa điện tử buồng cá nhân mở bằng thẻ RFID sinh viên hoặc nhận diện khuôn mặt (FaceID Check-in) tại cổng ký túc xá.
    \end{itemize}
    \item \textbf{Tích hợp AI \& Trợ lý Ảo Hỗ trợ Sinh viên:} Ứng dụng mô hình ngôn ngữ lớn (LLM) để xây dựng Chatbot hỗ trợ tự động giải đáp nội quy KTX, tiếp nhận thắc mắc và tự động phân loại mức độ khẩn cấp của các phiếu báo hỏng cơ sở vật chất.
\end{enumerate}



% ==============================================================================
% TÀI LIỆU THAM KHẢO
% ==============================================================================
\newpage
\addcontentsline{toc}{chapter}{TÀI LIỆU THAM KHẢO}
\bibliographystyle{IEEEtran}
\bibliography{references}

\end{document}

